\begin{abstract}
Differentiable opaque meshes, as introduced by \MS, render in any engine, but a trained scene occupies hundreds of megabytes and no compression method exists for them. We show that they cannot be compressed the way ordinary meshes or Gaussian splats are: because every pixel is covered by a single opaque triangle, simplifying a trained mesh loses 1.9\dB on average against training a mesh of that size from the start. The size of the file therefore has to be decided while the mesh is trained, and the one existing handle, a cap on the number of vertices, fixes neither the vertex count that survives training (0.68 to 1.06 of the cap) nor the bytes each vertex costs. We introduce \BM, a complete compression scheme spanning training, encoding and decoding. During training it runs the real encoder on the current mesh, converts a user-specified file size into a vertex count, grows the mesh toward that count by splitting the triangles with the largest rendering error, and trims the excess by edge collapse while optimization continues; a dedicated codec with a standalone decoder then writes one compact, engine-ready file. In one training run and without size statistics from other scenes, the file lands within 10\% of the requested size in 83 of 89 runs on four datasets. At 20 to 5~MB the files are 31--122$\times$ smaller than the trained model (7--27$\times$ smaller than its quantized, zstd-compressed PLY) at 0.3--1.3\dB lower PSNR and 0.09--0.19 higher LPIPS.
\end{abstract}
